% 1. Uvod
\section{Uvod}
\label{sec:uvod}

Ocenjivanje retko predstavlja najzanimljiviji deo nastave, ali oblikuje način na koji studenti uče.
Pravovremena i konkretna povratna informacija pomaže studentu da ispravi nesporazum pre nego što se ustali;
kada kasni ili je neodređena, provera znanja postaje formalnost.
Na velikim univerzitetskim kursevima nastavnici i saradnici moraju da ocene stotine odgovora otvorenog tipa u kratkim rokovima.
Kratka objašnjenja i mini-eseji, za razliku od pitanja sa ponuđenim odgovorima, ne mogu se proveriti prostim uparivanjem obrazaca.
Oni zahtevaju tumačenje: prepoznavanje delimično ispravnog rezonovanja,
razlikovanje jezičkih nepreciznosti od konceptualnih praznina i odluku o tome koliko neki propust treba da utiče na ocenu.
Takvo ocenjivanje je skupo, teško primenljivo na velike grupe studenata i sklono nedoslednosti:
čak i uz detaljna uputstva, dva ocenjivača mogu isti odgovor protumačiti različito,
naročito kada je odgovor dvosmislen, pisan u žurbi ili strukturno nejasan~\cite{brookhart2016century,jonsson2007use,stemler2004comparison}.
Nastavnika to opterećenje odvraća od pitanja zasnovanih na objašnjenjima,
a student povratnu informaciju dobija prekasno da bi mu koristila, ili je ne dobija uopšte.

Ranije generacije sistema za automatsku proveru znanja funkcionisale su samo kada su odgovori bili strogo ograničeni:
testovi sa ponuđenim odgovorima,
pitanja dopunjavanja ili kratki odgovori sa predvidljivom formulacijom~\cite{brew2013automated,ramesh2022automated,shermis2003automated}.
Napredak u obradi prirodnog jezika proširio je ove mogućnosti, ali su se mnogi pristupi i dalje oslanjali na ručno projektovanje atributa,
uske domene ili velike označene skupove podataka~\cite{banko2001scaling,collobert2008unified,sebastiani2002machine}.
Savremeni veliki jezički modeli (engl. \emph{large language models}, LLM) mogu da analiziraju i tumače odgovore formulisane slobodnim tekstom.
Time je postalo izvodljivo projektovati alate koji ocenjuju odgovore otvorenog tipa,
obrazlažu ocenu prirodnim jezikom i ukazuju na ono što u odgovoru nedostaje,
a ne samo na ono što je pogrešno~\cite{kasneci2023chatgpt,zhai2023chatgpt}.

\subsection{Motivacija}
\label{sec:motivacija}

Cilj ovog rada je ispitivanje mogućnosti pouzdane primene velikih jezičkih modela kao podrške u ocenjivanju.
Ručno ocenjivanje troši veliki deo nastavničkog vremena, pa povratna informacija kasni,
a provere znanja se organizuju ređe nego što bi bilo korisno.
Automatizovan sistem za ocenjivanje menja tu ravnotežu.
Provere znanja mogu biti češće, pa student kontinuirano prati napredak i prepoznaje oblasti koje zahtevaju dodatni rad.
Analiza odgovora stiže dok je gradivo još sveže.
Vreme koje bi otišlo na rutinski deo ocenjivanja nastavnik usmerava na individualnu podršku studentima.

Upotreba velikih jezičkih modela u ocenjivanju istovremeno otvara praktična i pedagoška pitanja.
Ovaj rad razmatra četiri:
\begin{enumerate}
    \item Koliko su ocene velikih jezičkih modela usklađene sa nastavničkim i koliko su dosledne kroz različite predmete i tipove pitanja?
    \item Mogu li nastavnici kalibrisati strogost ocenjivanja prema očekivanjima konkretnog kursa?
    \item Šta se dešava kada potpun referentni odgovor nastavnika nije dostupan?
    \item Kako sistem za automatizovano ocenjivanje projektovati tako da nastavnik zadrži kontrolu nad konačnom ocenom?
\end{enumerate}
Na ova pitanja rad se vraća u zaključku (poglavlje~\ref{sec:zakljucak}).

\subsection{Sistem EduGrader i doprinosi rada}

U ovom radu predstavljen je EduGrader, platforma za ocenjivanje pitanja otvorenog tipa u visokom obrazovanju, pod nadzorom nastavnika.
EduGrader objedinjuje više velikih jezičkih modela,
GPT-4o,\footnote{\url{https://platform.openai.com/docs/models/gpt-4o}}
GPT-4o-mini,\footnote{\url{https://platform.openai.com/docs/models/gpt-4o-mini}}
DeepSeek-Chat,\footnote{\url{https://api-docs.deepseek.com/}} i dva modela za rezonovanje (engl. \emph{reasoning models}),
GPT-5\footnote{\url{https://platform.openai.com/docs/models/gpt-5}}
i DeepSeek-Reasoner,\footnote{\url{https://api-docs.deepseek.com/guides/reasoning\_model}}
u jedinstven tok ocenjivanja u kome poslednju reč ima nastavnik.
Sistem podržava tri konfiguracije strogosti ocenjivanja: blagu, neutralnu i strogu.
Podela odgovara nastavnoj praksi, u kojoj neki kursevi nagrađuju delimično ispravno rezonovanje,
dok drugi zahtevaju preciznu terminologiju i potpuna rešenja.

EduGrader nudi dva komplementarna režima ocenjivanja:
\begin{itemize}
    \item \textbf{Režim sa zadatim referentnim odgovorom} (engl. \emph{provided-reference-answer}, PRA),
    u kome se studentski odgovori porede sa referentnim odgovorima koje je pripremio nastavnik.
    \item \textbf{Režim sa generisanim referentnim odgovorom} (engl. \emph{generated-reference-answer}, GRA),
    u kome sistem najpre generiše referentni odgovor na osnovu nastavnih materijala,
    beležaka sa predavanja, skripti i udžbenika,
    a zatim studentske odgovore ocenjuje u odnosu na tako generisani standard.
\end{itemize}

Izlaz u oba režima je numerička ocena praćena obrazloženjem koje ukazuje na jake strane odgovora, koncepte koji nedostaju i praznine u rezonovanju.
Takva povratna informacija studentu je upotrebljiva, a nastavniku proverljiva, za razliku od ocene koja dolazi iz „crne kutije''.

Model ocenu izražava na skali od 0 do 100, jer finija granulacija olakšava razlikovanje bliskih nivoa kvaliteta odgovora.
Dobijena vrednost se zatim preslikava na skalu od 0 do 10, na kojoj ocenjuju nastavnici (detalji u odeljku~\ref{sec:pipeline}).
Sve ocene u ovom radu, i sistemske i nastavničke, prikazane su na skali 0--10, osim tamo gde je naznačeno drugačije.

Svi eksperimenti sprovedeni su na srpskom jeziku,
koji spada u jezike sa ograničenim resursima i znatno je manje zastupljen u podacima za obučavanje savremenih modela,
čime je ispitivanje prošireno izvan dominantnog engleskog jezičkog konteksta.
Sistem je evaluiran na 583 autentična studentska odgovora prikupljena na dva univerziteta i pet kurseva,
od kratkih činjeničkih odgovora i otvorenih tekstualnih objašnjenja do odgovora u obliku koda.
Rezultati pokazuju da PRA režim daje stabilnije ishode od GRA režima i da modeli za rezonovanje daju ocene najbliže nastavničkim.
U najboljoj konfiguraciji (GPT-5, blaga strogost) EduGrader postiže Pirsonov koeficijent korelacije (engl.
\emph{Pearson correlation coefficient}, PCC) od 0{,}89 u odnosu na ljudsko ocenjivanje, a u neutralnoj konfiguraciji 0{,}88.

EduGrader je zamišljen kao alat za podršku ocenjivanju, a ne kao zamena za nastavnika.
Namenjen je smanjenju repetitivnog dela ocenjivanja,
ali ljudski nadzor ostaje neophodan za dvosmislena pitanja, očekivanja specifična za disciplinu i granične slučajeve.
Zato je sistem projektovan po principu čoveka u petlji (engl.
\emph{human-in-the-loop}) i objavljen kao softver otvorenog koda\footnote{\url{https://github.com/toswe/Master}},
kako bi bio proverljiv i prilagodljiv drugim institucijama.
Postupak ponovnog izvođenja svih prikazanih rezultata opisan je u odeljku~\ref{sec:reprodukcija}.

Doprinosi rada su sledeći:
\begin{enumerate}
    \item \textbf{Veb aplikacija EduGrader} (poglavlje~\ref{sec:aplikacija}):
    nastavnički i studentski modul, serverski modul za ocenjivanje koji čuva prompt, odgovor modela i izdvojenu ocenu za svaki ocenjeni odgovor,
    i integracioni sloj koji više provajdera jezičkih modela pokriva jednim interfejsom.
    \item \textbf{Sistem za automatizovanu evaluaciju modela} (poglavlje~\ref{sec:pipeline}):
    reproducibilan \emph{pipeline} koji nad ručno pripremljenim ispitnim podacima pokreće ocenjivanje
    za proizvoljnu kombinaciju modela, konfiguracije strogosti i režima i računa metrike korelacije i slaganja sa ocenama nastavnika.
    \item \textbf{Skup podataka na srpskom jeziku}: 583 autentična, pseudonimizovana studentska odgovora sa pet kurseva na dve institucije,
    sa nastavničkim ocenama i referentnim odgovorima (odeljci~\ref{sec:podaci} i~\ref{sec:etika}).
    \item \textbf{Sistematska evaluacija}: pet modela u tri konfiguracije strogosti i dva režima ocenjivanja,
    dopunjena eksperimentima sa nastavnim materijalom u promptu i sa temperaturom (poglavlje~\ref{sec:rezultati}).
    \item \textbf{Empirijski nalazi}: referentni odgovor nastavnika dosledno poboljšava i korelaciju i slaganje sa ljudskim ocenjivačima,
    modeli za rezonovanje su im najbliži, a stroga konfiguracija po pravilu pogoršava oba;
    najbolja konfiguracija dostiže PCC od 0{,}89.
\end{enumerate}

\subsection{Struktura rada}

Poglavlje~\ref{sec:pregled} daje pregled srodnih istraživanja primene veštačke inteligencije u ocenjivanju
i pozicionira ovaj rad u odnosu na njih.
Poglavlje~\ref{sec:aplikacija} opisuje veb aplikaciju EduGrader i sloj za integraciju jezičkih modela,
a poglavlje~\ref{sec:pipeline} sistem za automatizovanu evaluaciju modela kojim su sprovedeni eksperimenti.
Poglavlje~\ref{sec:metodologija} predstavlja metodologiju evaluacije: metrike, skupove podataka, zaštitu podataka i konfiguracije modela.
Poglavlje~\ref{sec:rezultati} donosi rezultate, uključujući poređenje PRA i GRA režima kroz modele i konfiguracije strogosti,
uticaj konfiguracije strogosti i korelacije po pojedinačnim kursevima.
U poglavlju~\ref{sec:diskusija} rezultati su upoređeni sa prethodnim istraživanjima i razmotreni su pravci budućeg rada,
dok poglavlje~\ref{sec:zakljucak} zaključuje rad.
Dodatak~\ref{app:promptovi} sadrži šablone promptova (engl. \emph{prompt}) i reprezentativne primere ocenjivanja.
